
\chapter*{Abstract}
\label{chap:abstract}
\thispagestyle{empty}

The increasing integration of \gls{ai}-based systems into \gls{df} workflows
is changing how large collections of digital images are classified, retrieved,
and prioritized for human examination. Although these systems may reduce
analyst workload and accelerate the identification of potentially relevant
content, their operational value depends on whether their behavior remains
stable when inputs are degraded, distributionally different, or deliberately
manipulated.

This thesis evaluates the operational robustness of transparent proxy models
and selected black-box software configurations used for automated image
classification and semantic retrieval in forensic contexts. Rather than
proposing a new classifier or ranking commercial products, the study examines
how heterogeneous systems behave under clean, \gls{ood}, adversarial, and
anti-forensic conditions. The evaluation is organized through \gls{fairlab}, a
task-specific methodological protocol connecting human-validated dataset
construction, deterministic partitioning, transparent proxy-model evaluation,
perturbation generation, blind-bundle preparation, black-box output
normalization, quantitative analysis, qualitative explainability, and integrity
controls. Errors are interpreted in terms of missed-target risk, human-review
workload, directional transitions, configuration sensitivity, and
source-to-output traceability.

The frozen source population contains 1\,500 images: 500
\texttt{weapon} images, 500 \texttt{non\_weapon} images, and a separate
500-image \gls{ood} branch. The \gls{ood} population is not treated as a third
supervised class. It is used to characterize the assignments produced when
closed-set systems process images that do not conform to the nominal binary
task.

The 1\,000-image binary subset provides the common basis for generating 5\,000
adversarial or adversarial-style artifacts and 5\,000 anti-forensic artifacts.
Together with the 1\,000 clean binary inputs and 500 clean \gls{ood} inputs,
these derived populations form an 11\,500-item forensic evaluation bundle. The
10\,000 manipulated artifacts remain source--condition variants of the same
1\,000 binary images and do not represent independent underlying cases.

The transparent level comprises \gls{efficientnetb0}, \gls{resnet18}, and a
\gls{clip}-based proxy. The black-box level comprises four software systems
evaluated through six frozen configurations:
\gls{magnetaxiomtool}/\gls{magnetai},
\gls{griffeye}/\gls{t3kcore},
\gls{excirefoto2025} at \texttt{D20}, \texttt{D50}, and \texttt{D80}, and
\gls{cellebriteinseyets}. Their tags, bookmarks, exported labels, and
semantic-retrieval results are interpreted through documented normalization
rules without assuming architectural or semantic equivalence.

Four model-dependent attacks---\gls{fgsm}, \gls{onepixel}, \gls{sigmazero},
and \gls{superdeepfool}---were generated against fold-specific
\gls{efficientnetb0} checkpoints. Their evaluation on \gls{efficientnetb0}
constitutes direct-target assessment, whereas their application to the other
proxies and black-box configurations measures empirical transfer. The
model-agnostic \gls{colorshift} condition and five anti-forensic transformations
are interpreted separately. The analysis combines conventional performance
metrics with directional error transitions, \gls{ood} assignment behavior,
architecture-specific maximum predicted-class probabilities, and five
qualitative \gls{ig} case studies. Probability and explainability outputs are
used only as model-specific diagnostic evidence.

The results show that strong clean performance does not imply uniform
operational robustness. Clean proxy-model accuracies are 0.978 for
\gls{efficientnetb0}, 0.964 for \gls{resnet18}, and 0.927 for the
\gls{clip}-based proxy. Under direct-target evaluation,
\gls{efficientnetb0} accuracy decreases to 0.575 under \gls{fgsm} and to 0.015
under \gls{sigmazero}. The same artifacts produce more limited and
heterogeneous effects on the non-target systems, demonstrating that empirical
transfer must be measured rather than assumed. The anti-forensic conditions
generally produce smaller aggregate variations than the strongest direct-target
attacks but still introduce operationally relevant sample-level errors.

The separate \gls{ood} analysis further shows that out-of-task inputs may still
be assigned to the firearm-oriented operational state. At the proxy level, the
fold-aggregated weapon assignment rates are 0.3696 for \gls{efficientnetb0},
0.4376 for \gls{resnet18}, and 0.8356 for the \gls{clip}-based proxy. These
values represent closed-set assignment behavior rather than supervised
correctness.

The black-box evaluation produces 69\,000 normalized decisions and reveals
distinct configuration-dependent operating profiles. No evaluated
configuration simultaneously minimizes missed detections, false-positive review
workload, and \gls{ood} weapon assignments. In particular, changes to the
\gls{excirefoto2025} \texttt{Distance limit} materially alter recall,
selectivity, and \gls{ood} behavior. The results therefore characterize the
frozen software versions, enabled functions, configurations, observable output
semantics, and normalization mappings.

The thesis concludes that \gls{ai}-based image-analysis systems may provide
useful support for forensic triage but must not be treated as autonomous
mechanisms for evidentiary decision-making. Their defensible use requires
task-specific validation, explicit threat modeling, documented versions and
configurations, separate analysis of \gls{ood} behavior, source-to-output
traceability, and continued expert verification. \gls{fairlab} contributes a
unified and auditable procedure for making these requirements experimentally
observable across transparent proxy models and selected black-box software
configurations evaluated on the same frozen experimental basis.

\vfill

\noindent
\begin{minipage}{\textwidth}
\small
\raggedright
\textbf{Keywords:} Digital Forensics; Artificial Intelligence; AI-based forensic
tools; image classification; operational robustness; adversarial attacks;
anti-forensic techniques; out-of-distribution analysis; explainability.
\end{minipage}
